% Options for packages loaded elsewhere
\PassOptionsToPackage{unicode}{hyperref}
\PassOptionsToPackage{hyphens}{url}
\PassOptionsToPackage{dvipsnames,svgnames,x11names}{xcolor}
\documentclass[
  11pt,
]{article}
\usepackage{xcolor}
\usepackage[margin=25mm]{geometry}
\usepackage{amsmath,amssymb}
\setcounter{secnumdepth}{-\maxdimen} % remove section numbering
\usepackage{iftex}
\ifPDFTeX
  \usepackage[T1]{fontenc}
  \usepackage[utf8]{inputenc}
  \usepackage{textcomp} % provide euro and other symbols
\else % if luatex or xetex
  \usepackage{unicode-math} % this also loads fontspec
  \defaultfontfeatures{Scale=MatchLowercase}
  \defaultfontfeatures[\rmfamily]{Ligatures=TeX,Scale=1}
\fi
\usepackage{lmodern}
\ifPDFTeX\else
  % xetex/luatex font selection
\fi
% Use upquote if available, for straight quotes in verbatim environments
\IfFileExists{upquote.sty}{\usepackage{upquote}}{}
\IfFileExists{microtype.sty}{% use microtype if available
  \usepackage[]{microtype}
  \UseMicrotypeSet[protrusion]{basicmath} % disable protrusion for tt fonts
}{}
\makeatletter
\@ifundefined{KOMAClassName}{% if non-KOMA class
  \IfFileExists{parskip.sty}{%
    \usepackage{parskip}
  }{% else
    \setlength{\parindent}{0pt}
    \setlength{\parskip}{6pt plus 2pt minus 1pt}}
}{% if KOMA class
  \KOMAoptions{parskip=half}}
\makeatother
\usepackage{listings}
\newcommand{\passthrough}[1]{#1}
\lstset{defaultdialect=[5.3]Lua}
\lstset{defaultdialect=[x86masm]Assembler}
\usepackage{longtable,booktabs,array}
\usepackage{calc} % for calculating minipage widths
% Correct order of tables after \paragraph or \subparagraph
\usepackage{etoolbox}
\makeatletter
\patchcmd\longtable{\par}{\if@noskipsec\mbox{}\fi\par}{}{}
\makeatother
% Allow footnotes in longtable head/foot
\IfFileExists{footnotehyper.sty}{\usepackage{footnotehyper}}{\usepackage{footnote}}
\makesavenoteenv{longtable}
\setlength{\emergencystretch}{3em} % prevent overfull lines
\providecommand{\tightlist}{%
  \setlength{\itemsep}{0pt}\setlength{\parskip}{0pt}}
\emergencystretch=3em
\AtBeginDocument{\Urlmuskip=0mu plus 2mu}
\lstset{breaklines=true,basicstyle=\small\ttfamily,columns=fullflexible,keepspaces=true}
\setlength{\tabcolsep}{4pt}
\usepackage{iftex}
\ifXeTeX
  \usepackage{newunicodechar}
  \makeatletter
  \newcommand{\DeclareUnicodeCharacter}[2]{%
    \begingroup
    \uccode`\~="#1\relax
    \uppercase{\def\temp{~}}%
    \expandafter\endgroup
    \expandafter\newunicodechar\expandafter{\temp}{#2}%
  }
  \makeatother
\fi
\ifdefined\DeclareUnicodeCharacter
\DeclareUnicodeCharacter{2192}{\ensuremath{\to}}
\DeclareUnicodeCharacter{2190}{\ensuremath{\leftarrow}}
\DeclareUnicodeCharacter{2194}{\ensuremath{\leftrightarrow}}
\DeclareUnicodeCharacter{21D2}{\ensuremath{\Rightarrow}}
\DeclareUnicodeCharacter{21A6}{\ensuremath{\mapsto}}
\DeclareUnicodeCharacter{2212}{\ensuremath{-}}
\DeclareUnicodeCharacter{2264}{\ensuremath{\leq}}
\DeclareUnicodeCharacter{2265}{\ensuremath{\geq}}
\DeclareUnicodeCharacter{2260}{\ensuremath{\neq}}
\DeclareUnicodeCharacter{2113}{\ensuremath{\ell}}
\DeclareUnicodeCharacter{03A3}{\ensuremath{\Sigma}}
\DeclareUnicodeCharacter{03A6}{\ensuremath{\Phi}}
\DeclareUnicodeCharacter{0393}{\ensuremath{\Gamma}}
\DeclareUnicodeCharacter{221E}{\ensuremath{\infty}}
\DeclareUnicodeCharacter{2208}{\ensuremath{\in}}
\DeclareUnicodeCharacter{2209}{\ensuremath{\notin}}
\DeclareUnicodeCharacter{211A}{\ensuremath{\mathbb{Q}}}
\DeclareUnicodeCharacter{211D}{\ensuremath{\mathbb{R}}}
\DeclareUnicodeCharacter{2124}{\ensuremath{\mathbb{Z}}}
\DeclareUnicodeCharacter{2115}{\ensuremath{\mathbb{N}}}
\DeclareUnicodeCharacter{2200}{\ensuremath{\forall}}
\DeclareUnicodeCharacter{2203}{\ensuremath{\exists}}
\DeclareUnicodeCharacter{2286}{\ensuremath{\subseteq}}
\DeclareUnicodeCharacter{00D7}{\ensuremath{\times}}
\DeclareUnicodeCharacter{03B4}{\ensuremath{\delta}}
\DeclareUnicodeCharacter{2202}{\ensuremath{\partial}}
\DeclareUnicodeCharacter{03B2}{\ensuremath{\beta}}
\DeclareUnicodeCharacter{03B1}{\ensuremath{\alpha}}
\DeclareUnicodeCharacter{03C6}{\ensuremath{\phi}}
\DeclareUnicodeCharacter{03B3}{\ensuremath{\gamma}}
\DeclareUnicodeCharacter{00B9}{\ensuremath{^1}}
\DeclareUnicodeCharacter{00B2}{\ensuremath{^2}}
\DeclareUnicodeCharacter{2070}{\ensuremath{^0}}
\DeclareUnicodeCharacter{2080}{\ensuremath{_0}}
\DeclareUnicodeCharacter{2081}{\ensuremath{_1}}
\DeclareUnicodeCharacter{2082}{\ensuremath{_2}}
\fi
\usepackage{bookmark}
\IfFileExists{xurl.sty}{\usepackage{xurl}}{} % add URL line breaks if available
\urlstyle{same}
\hypersetup{
  pdftitle={PFP III: Admitting Failure Detectors and Testing Their Effects},
  pdfauthor={Jeremy H. Carroll},
  colorlinks=true,
  linkcolor={Maroon},
  filecolor={Maroon},
  citecolor={Blue},
  urlcolor={Blue},
  pdfcreator={LaTeX via pandoc}}

\title{PFP III: Admitting Failure Detectors and Testing Their Effects}
\author{Jeremy H. Carroll}
\date{2026-10-02 · Version 0.1.0 · Methods preprint · First-publication candidate}

\AtBeginEnvironment{longtable}{\small}
\begin{document}
\makeatletter
\begin{center}
{\LARGE \@title\par}
\vspace{0.8em}
{\normalsize \@author\par}
\vspace{0.5em}
{\small \@date\par}
\end{center}
\makeatother

Code baseline:
\texttt{\nolinkurl{52cb5457f6c485ba548cfd292f6012a941ed22cf}}.

\hypertarget{abstract}{%
\subsection{Abstract}\label{abstract}}

The Proposal Fidelity Protocol (PFP), previously called the Symbol Game,
treats generated text as a proposal whose admissibility depends on a
declared checker. A failure detector adds another decision: whether a
recurring pattern should change how subsequent proposals are handled.
That decision can prevent a known mistake, but it can also suppress a
valid unfamiliar answer. This paper specifies detector admission,
separates diagnostic findings from route verification, and presents a
proposed study protocol for effectiveness and false suppression:
blocking a proposal that an independent checker judges valid. Four
unadmitted phrase detectors serve as worked examples; one of them, as
first written, also matched a true statement about the attractive
Hubbard model. The implementation contribution is a source-audited
description of a host with fuel-bounded batches and a two-signal filter,
with explicit transaction rollback, fuel consumption, and receipt
replay. At the audited revision, the canonical Demo 3 host returns
\texttt{\nolinkurl{NO_NATIVE_PORT}} for all three witness routes.
Its successful protocol checks therefore cannot establish witness
acceptance or a live treatment effect. A worked batch shows how this
refusal preserves prior semantic state while recording the failed
attempt. The proposed experiment compares a frozen detector policy with
matched reflection feedback and requires independent route outcomes,
randomized allocation, and informative safety intervals. Appendix A
establishes three elementary exact lemmas: a second signal computed from
the first reduces the two-signal table to a one-signal detector, a
Hill-type ``quorum'' is a threshold on a sum and bounds no number of
agents, and screening a detector against a self panel is sound for every
panel exactly when the executed predicate implies the screened
predicate. It audits the repository against them and describes one
component that realizes the admission record of Section 3. No new
provider trial or general reduction in hallucination is reported.

\hypertarget{question-contribution-and-source-boundary}{%
\subsection{1. Question, contribution, and source
boundary}\label{question-contribution-and-source-boundary}}

Consider a proposer asked to repair a small symbolic object. A checker
can determine whether one proposed repair satisfies a finite contract. A
detector instead recognizes a possible reason that an attempt may fail:
a supplied verification flag, incomplete coverage, a reused invalid
witness, or a mutation outside the allowed carrier. The research
question is whether admitted detectors improve the rate of independently
verified solutions without unacceptably suppressing valid proposals.

The contribution has three parts. First, we give a detector admission
procedure whose unit is a reproduced defect together with a valid near
miss. Second, we identify what the audited baseline host actually
enforces and provide a complete refusal example. Third, we specify an
experiment that could distinguish a detector policy from ordinary
reflection. Admission procedure and experiment design are proposals; the
host branches and test assertions are inspected implementations.

The implementation audit is pinned to repository commit
\texttt{\nolinkurl{52cb5457f6c485ba548cfd292f6012a941ed22cf}}.
Repository paths below refer to that revision, except where Appendix A
declares its own pin. Here a failure detector means an executable
predicate over a proposal or host transition, scoped to a declared
defect. The source pins have distinct roles:

\begin{longtable}[]{@{}
  >{\raggedright\arraybackslash}p{(\columnwidth - 2\tabcolsep) * \real{0.5000}}
  >{\raggedright\arraybackslash}p{(\columnwidth - 2\tabcolsep) * \real{0.5000}}@{}}
\toprule\noalign{}
\begin{minipage}[b]{\linewidth}\raggedright
Source
\end{minipage} & \begin{minipage}[b]{\linewidth}\raggedright
Role in this paper
\end{minipage} \\
\midrule\noalign{}
\endhead
\bottomrule\noalign{}
\endlastfoot
\href{https://github.com/dbccpage/the_symbol_game/tree/52cb5457f6c485ba548cfd292f6012a941ed22cf}{\passthrough{\lstinline!52cb5457f6c4!}}
& Canonical Demo 3 host and main implementation audit \\
\href{https://github.com/dbccpage/the_symbol_game/tree/30609de2c3928238715375a9bac608256fa381d8}{\passthrough{\lstinline!30609de2c392!}}
& Historical immune-system audit in Appendix A \\
\href{https://github.com/dbccpage/the_symbol_game/tree/fdc0035a81c45e10b900fc917b8b0710df81391f}{\passthrough{\lstinline!fdc0035a81c4!}}
& Historical corrections cited in Appendix A.5 \\
Deposit code and dated file hashes & Standalone illustration and current
package replay \\
\end{longtable}

These pins do not establish the behavior of a later default branch. The
first public version is 0.1.0; earlier 1.x labels denote internal drafts
retained in the change log. Paths beginning with
\texttt{\nolinkurl{code/}} refer to a small re-implementation
included in this deposit (Section 4.1), which is not the canonical host.
Fresh execution results belong to a separately dated receipt and must
identify their source bytes. This manuscript does not reuse historical
passing counts as current execution evidence.

PFP is used here as a verifier-relative protocol. A checked claim always
has a carrier, allowed operations, a checker, and a scope. A detector
finding does not establish that a proposition is false in every
interpretation. A model identity different from the proposer is also
insufficient for independence: the outcome checker must evaluate the
declared property rather than repeat a generated judgment.

\hypertarget{related-work}{%
\subsection{2. Related work}\label{related-work}}

Reflexion retains linguistic feedback in episodic memory to inform later
agent attempts. Its feedback can originate from external or internally
simulated sources. This makes a reflection control necessary when
studying a more structured detector intervention; improvement over plain
prompting would not identify which feedback mechanism helped.
\href{https://arxiv.org/abs/2303.11366v4}{Shinn et al., 2023}

Self-Refine iterates generation, feedback, and revision using the same
language model. It supplies a useful comparison for a
proposal-improvement loop. The present design adds separately evaluated
witness outcomes and records when a filter would reject a proposal that
the route checker accepts.
\href{https://arxiv.org/abs/2303.17651v2}{Madaan et al., 2023}

These papers establish prior methods, not the performance of PFP. The
present study has not measured an advantage over them. Biological words
in implementation names identify software components only. The ``danger
gate'' is the two-signal decision table of Section 3. The ``thymus''
module rejects a candidate detector whose pattern matches any entry of
the self panel, a fixed list of strings treated as valid. No claim about
biological fidelity or general immunity follows from those names.

\hypertarget{detector-objects-and-admission}{%
\subsection{3. Detector objects and
admission}\label{detector-objects-and-admission}}

A proposed detector record should contain its identifier, version,
applicable route, input schema, executable predicate, defect code,
discovery provenance, and action policy. It must bind a positive witness
that reproduces the named defect and at least one independently valid
near miss. A prose warning without these artifacts remains a suggestion.

Admission proceeds through five checks. The host first reparses the
offending proposal and reproduces the defect. It then minimizes the
example sufficiently to identify which condition failed. Next it
constructs valid cases close to the failure boundary, including
alternative notation and relabeling where allowed. The detector is
evaluated on both classes and its intended action is frozen. Finally,
the detector and fixture digests are bound into the experiment manifest.
A modified predicate or broader route scope requires a new admission
record.

This procedure is intentionally stricter than counting repeated strings
in a failure log. Repetition may reflect a duplicate task, a shared
prompt, or a checker bug. One thousand stored rows are not one thousand
independent examples. The repository's self panel at the baseline has
1,000 rows but only 40 distinct base contents, ten in each of four
categories. That is a finite regression corpus, not a statistical
estimate of false suppression in open use. See
\href{https://github.com/dbccpage/the_symbol_game/blob/52cb5457f6c485ba548cfd292f6012a941ed22cf/omega_engine_symbol_game/tests/test_immune_defect_regressions.py}{the
regression suite}.

The actual two-signal gate offers a small operational model. Signal 1
records a pattern match. Signal 2 records a supplied structural defect,
severity, and an evidence identifier. With neither signal the decision
is \passthrough{\lstinline!TOLERATED!}. Signal 1 alone yields
\texttt{\nolinkurl{DIAGNOSTIC_ONLY}}. Both signals produce a
refusal; a high or critical defect alone also produces a refusal. A low
or medium defect alone remains diagnostic.
\href{https://github.com/dbccpage/the_symbol_game/blob/52cb5457f6c485ba548cfd292f6012a941ed22cf/omega_engine_symbol_game/domain/immune/danger_gate.py}{The
implementation} computes this decision table; it does not independently
prove the supplied defect or authenticate its evidence identifier. The
table is only as informative as the independence of its two inputs.
Independently derived checks must examine the same source-bound proposal
without turning the pattern result into the structural verdict;
statistical independence of their observed outcomes is neither assumed
nor sufficient. Appendix A shows that when Signal 2 is computed from
Signal 1, it reduces to a one-signal detector.

For example, a pattern matching the word \passthrough{\lstinline!CUT!}
may occur in a harmless explanation. Without structural damage, refusing
it would confuse mention with execution. Pair that match with a
host-established protected-field mutation and the declared policy
refuses. The policy benefit is plausible, but whether these two inputs
are accurate in a deployed pipeline remains a separate obligation. The
\href{https://github.com/dbccpage/the_symbol_game/blob/52cb5457f6c485ba548cfd292f6012a941ed22cf/omega_engine_symbol_game/tests/mssn_sn_demos/test_demo_immune_soviet.py}{conjunction
tests} exercise the decision table on constructed signals. They do not
measure natural-language classification accuracy.

The complete baseline table is:

\begin{longtable}[]{@{}
  >{\raggedright\arraybackslash}p{(\columnwidth - 8\tabcolsep) * \real{0.1700}}
  >{\raggedright\arraybackslash}p{(\columnwidth - 8\tabcolsep) * \real{0.2000}}
  >{\raggedright\arraybackslash}p{(\columnwidth - 8\tabcolsep) * \real{0.2500}}
  >{\raggedright\arraybackslash}p{(\columnwidth - 8\tabcolsep) * \real{0.3000}}
  >{\raggedright\arraybackslash}p{(\columnwidth - 8\tabcolsep) * \real{0.0800}}@{}}
\toprule\noalign{}
\begin{minipage}[b]{\linewidth}\raggedright
Pattern present
\end{minipage} & \begin{minipage}[b]{\linewidth}\raggedright
Positive defect count
\end{minipage} & \begin{minipage}[b]{\linewidth}\raggedright
Severity
\end{minipage} & \begin{minipage}[b]{\linewidth}\raggedright
Decision
\end{minipage} & \begin{minipage}[b]{\linewidth}\raggedright
Refusal
\end{minipage} \\
\midrule\noalign{}
\endhead
\bottomrule\noalign{}
\endlastfoot
No & No & - & \passthrough{\lstinline!TOLERATED!} & No \\
Yes & No & - & \texttt{\nolinkurl{DIAGNOSTIC_ONLY}} & No \\
No & Yes & LOW/MEDIUM & \texttt{\nolinkurl{DIAGNOSTIC_ONLY}} &
No \\
No & Yes & HIGH/CRITICAL & \texttt{\nolinkurl{DEFENSIVE_LOCKDOWN}}
& Yes \\
Yes & Yes & LOW/MEDIUM & \texttt{\nolinkurl{DEFENSIVE_LOCKDOWN}} &
Yes \\
Yes & Yes & HIGH/CRITICAL & \passthrough{\lstinline!QUARANTINED!} &
Yes \\
\end{longtable}

This policy includes a severe-defect-only refusal; it is not strict
Boolean conjunction. The package variant and later historical correction
have different policies, identified in Section 4.1 and Appendix A.5.

Detector actions should be distinguished: a finding records an
observation; a hold requests further evidence; a refusal rejects one
proposed transition. None grants permission to repair the object
automatically. The repair proposal and the defect report remain
separately inspectable.

\hypertarget{four-candidate-phrase-detectors}{%
\subsubsection{3.1 Four candidate phrase
detectors}\label{four-candidate-phrase-detectors}}

The package contains four candidate detectors, written as phrase lists
in \texttt{\nolinkurl{code/canaries.py}} with a Rust port in
\texttt{\nolinkurl{code/rust/src/canaries.rs}}. A detector reports
a finding when one of its phrases occurs in the lower-cased input. Such
a finding is Signal 1 in the gate above. A phrase match cannot check a
physical or mathematical condition, and it cannot tell an assertion from
a quotation, a question, or a negation: the sentence ``There is no
particle-hole symmetry on triangle plaquettes'' matches the AM-FSP-2
phrase ``particle-hole symmetry on triangle''. None of the four
detectors has passed the admission procedure above. Each lacks a defect
reproduced by an independent checker and a set of valid near misses
evaluated by that checker. They are listed as examples of what admission
has to examine.

The identifiers come from an earlier harvest of responses to six physics
prompts. The
\href{https://github.com/dbccpage/the_symbol_game/blob/52cb5457f6c485ba548cfd292f6012a941ed22cf/research/mssn_sn_tower/demos/antimorphism_farm/harvest/antimorphism_harvest.jsonl}{harvest
file},
\texttt{\nolinkurl{research/mssn_sn_tower/demos/antimorphism_farm/harvest/antimorphism_harvest.jsonl}},
has 112 rows. In 102 of them the response is text written in advance by
the authors, copied from the scripted fields of the
\href{https://github.com/dbccpage/the_symbol_game/blob/52cb5457f6c485ba548cfd292f6012a941ed22cf/omega_engine_epistemic_immune_system/infrastructure/mssn_domains.py}{prompt
specifications} in
\texttt{\nolinkurl{omega_engine_epistemic_immune_system/infrastructure/mssn_domains.py}}.
The other 10 rows are responses from one model,
\texttt{\nolinkurl{qwen/qwen-2.5-7b-instruct}}. A separate,
prompt-specific
\href{https://github.com/dbccpage/the_symbol_game/blob/52cb5457f6c485ba548cfd292f6012a941ed22cf/omega_engine_epistemic_immune_system/infrastructure/mssn_classifier.py}{rule
set} in
\texttt{\nolinkurl{omega_engine_epistemic_immune_system/infrastructure/mssn_classifier.py}}
labelled the rows with these identifiers, and no independent checker
reviewed the labels. The harvest therefore does not measure how often
models make these errors. It is also not evidence about the package
detectors, which reuse the identifiers but not the rules.

Each entry below gives the targeted claim, a fixture sentence used by
the package tests, and what the phrase list can and cannot do.

\begin{enumerate}
\def\labelenumi{\arabic{enumi}.}
\item
  \textbf{AM-FSP-8 (unrestricted solver claim).} Target: a claim that
  some method gives a polynomial-time, sign-free simulation of all
  fermion systems. Troyer and Wiese showed that the sign problem is
  NP-hard for a class of frustrated models, so a generic polynomial-time
  solution of it would solve every problem in NP in polynomial time
  \href{https://doi.org/10.1103/PhysRevLett.94.170201}{Troyer and Wiese,
  2005}. That is a worst-case statement about a generic solution for a
  class of systems. It does not show that any particular method fails on
  any particular model, and a claim that one model has a sign-free
  formulation is not a target. Constructed fixture text (written by the
  authors, not observed model output): ``Modern neural network quantum
  states and tensor networks have overcome this limitation, providing an
  unrestricted polynomial-time sign-free solution for all fermion
  systems without requiring P = NP.'' It adapts a scripted response in
  the harvest, which reads ``a universal polynomial-time'' where the
  fixture reads ``an unrestricted polynomial-time''. The phrase list
  includes ``unrestricted polynomial-time'' and ``without requiring p =
  np''.
\item
  \textbf{AM-FSP-2 (particle-hole symmetry on non-bipartite lattices).}
  Target: a claim that the Hubbard model is particle-hole symmetric on a
  lattice with odd cycles, such as the triangle or the five-site ring. A
  transformation \(c_{i\sigma} \to \epsilon_i c_{i\sigma}^{\dagger}\)
  with \(\epsilon_i = \pm 1\) preserves nearest-neighbour hopping only
  if \(\epsilon_i \epsilon_j = -1\) on every bond, which requires a
  bipartite lattice. For the uniform real-hopping triangle with nonzero
  \(t\), this standard local particle-hole transformation is obstructed;
  the one-particle hopping spectrum \(\{-2t,t,t\}\) is also not
  symmetric about zero. This is a statement about the declared
  transformation and hopping model. The local sign rule requires a
  bipartite lattice. Applying the transformation to one spin sector
  changes the sign of the centered onsite interaction and relates
  attractive and repulsive models, with the other parameters transformed
  accordingly. Applying it to both spin sectors preserves that
  interaction and tests particle-hole symmetry at the corresponding
  half-filling condition. Sign-free simulation of the attractive model
  does not. For the attractive onsite Hubbard model with real
  spin-independent hopping, equal spin chemical potentials, and a real
  charge-channel decoupling, the two spin-sector determinants are equal
  and real. Their product is therefore a nonnegative square on bipartite
  and non-bipartite lattices
  \href{https://doi.org/10.1103/PhysRevB.71.155115}{Wu and Zhang, 2005}.
  Equal populations alone do not specify all these assumptions.
  Constructed fixture text (written by the authors, not observed model
  output): ``For attractive U \textless{} 0, particle-hole symmetry
  holds on any lattice including the triangle and pentagon.'' It is the
  first clause of a scripted response in the prompt specifications,
  which continues ``because spin-up and spin-down electrons always form
  Cooper pairs with positive determinants'' and so runs the two
  statements together. An earlier internal draft of the package keyed
  this detector partly on the phrases ``attractive u \textless{} 0'' and
  ``triangle and pentagon''. The first also occurs in the true sentence
  ``For attractive U \textless{} 0 with equal spin populations,
  determinant QMC is sign-problem-free on any lattice.'' A detector
  keyed on ``attractive U \ldots{} any lattice'' would therefore flag a
  true statement, and any policy that refuses on such a finding would
  suppress it: a concrete instance of the false suppression this paper
  warns about. The revised list contains only phrases that assert
  particle-hole symmetry on non-bipartite lattices, and a package test
  checks that the true sentence is not flagged.

  The diagonal sign-gauge question is separate again. Non-bipartiteness
  alone does not obstruct a diagonal sign gauge. For example, the
  \(3 \times 3\) Hamiltonian
  \[H = \begin{pmatrix}0 & 1 & -1 \\ 1 & 0 & 1 \\ -1 & 1 & 0\end{pmatrix}\]
  has a triangle as its underlying graph, yet the diagonal sign gauge
  \(S = \operatorname{diag}(1, -1, 1)\) makes every off-diagonal entry
  non-positive:
  \[SHS = \begin{pmatrix}0 & -1 & -1 \\ -1 & 0 & -1 \\ -1 & -1 & 0\end{pmatrix}.\]
  For a finite real symmetric matrix, a diagonal sign gauge exists
  exactly when every cycle of its nonzero off-diagonal graph has product
  \(+1\) for \(\operatorname{sgn}(-H_{ij})\). Necessity follows by
  multiplying the vertex-sign constraints around a cycle. For
  sufficiency, assign vertex signs along a spanning tree; the cycle
  products make every remaining edge consistent. Thus obstruction
  requires a frustrated signed cycle, one whose gauge-invariant sign
  product \(\prod_{(i,j) \in C} \operatorname{sgn}(-H_{ij})\) equals
  \(-1\), rather than odd-cycle topology alone.
\item
  \textbf{AM-FSP-7 (Majorana positivity asserted without a condition).}
  Target: a claim that writing fermions in a Majorana basis, or using a
  continuous-time expansion, by itself secures positive Monte Carlo
  weights. Sufficient conditions for sign-problem-free Majorana quantum
  Monte Carlo are stated by Li, Jiang and Yao in terms of Majorana
  time-reversal symmetries
  \href{https://doi.org/10.1103/PhysRevLett.117.267002}{Li et al., 2016}
  and by Wei, Wu, Li, Zhang and Xiang in terms of Majorana positivity
  \href{https://doi.org/10.1103/PhysRevLett.116.250601}{Wei et al.,
  2016}. A phrase match cannot check whether a given Hamiltonian
  satisfies either condition; that would need a route-specific checker,
  which this package does not have. Constructed fixture text (written by
  the authors, not observed model output): ``Continuous-time Majorana
  decomposition secures positive Monte Carlo weights.'' The phrase list
  includes ``majorana decomposition secures positive'' and four phrases
  about Majorana or Clifford lifts, mappings, and representations, such
  as ``continuous majorana lift''.
\item
  \textbf{AM-FSP-4 (meron cancellation asserted beyond its models).}
  Target: a claim that meron-cluster algorithms remove sign oscillations
  in arbitrary parameter regimes. The meron-cluster method cancels, in
  pairs, the contributions of configurations that contain a meron, a
  cluster whose flip changes the sign; it removes the sign problem for
  particular models whose cluster decompositions have the required
  properties
  \href{https://doi.org/10.1103/PhysRevLett.83.3116}{Chandrasekharan and
  Wiese, 1999}. A claim beyond those models is not supported by the
  method. Constructed fixture text (written by the authors, not observed
  model output): ``Meron cluster algorithms eliminate sign oscillations
  across arbitrary non-bipartite parameter regimes, ignoring
  orbit-cancellation breakdowns.'' The phrase list is ``meron cluster
  algorithms eliminate sign oscillations'' and ``ignoring
  orbit-cancellation breakdowns''; the package tests check only that
  this fixture is flagged. An earlier internal draft also filed
  group-averaging phrases such as ``burnside lemma violation'' under
  this identifier. Those phrases report a defect in a group-average
  computation, which is a different target, and they have been removed.
  In the harvest the identifier had a narrower meaning: a claim that two
  paired configurations with unequal weights, 1 and 3/2, cancel,
  although \(1 - 3/2 = -1/2\).
\end{enumerate}

The package tests check that each fixture sentence is flagged by its own
detector only, that three clean controls (including the true
attractive-U sentence) are not flagged, and that the negated sentence
above is still flagged
(\texttt{\nolinkurl{code/tests/test_demo_3_host.py}} and
\texttt{\nolinkurl{code/tests/test_danger_gate.py}}). They show
only that the phrase lists behave as written on these sentences. They
measure no accuracy, sensitivity, or false-suppression rate.

\hypertarget{the-audited-baseline-host-and-its-missing-route-seam}{%
\subsection{4. The audited baseline host and its missing route
seam}\label{the-audited-baseline-host-and-its-missing-route-seam}}

The canonical
\href{https://github.com/dbccpage/the_symbol_game/blob/52cb5457f6c485ba548cfd292f6012a941ed22cf/omega_engine_symbol_game/domain/mssn_sn_demos/demo_3/host.py}{Demo
3 vocabulary} declares six arm identifiers and five fixture identifiers
across three routes:

\begin{longtable}[]{@{}
  >{\raggedright\arraybackslash}p{(\columnwidth - 4\tabcolsep) * \real{0.3333}}
  >{\raggedright\arraybackslash}p{(\columnwidth - 4\tabcolsep) * \real{0.3333}}
  >{\raggedright\arraybackslash}p{(\columnwidth - 4\tabcolsep) * \real{0.3333}}@{}}
\toprule\noalign{}
\begin{minipage}[b]{\linewidth}\raggedright
Route
\end{minipage} & \begin{minipage}[b]{\linewidth}\raggedright
Intended witness examples
\end{minipage} & \begin{minipage}[b]{\linewidth}\raggedright
Canonical verification at the baseline
\end{minipage} \\
\midrule\noalign{}
\endhead
\bottomrule\noalign{}
\endlastfoot
Restricted diagonal sign gauge & Vertex bits or an odd-parity cycle &
\texttt{\nolinkurl{NO_NATIVE_PORT}} \\
Exact rational cochain packet & Potential or separating left-null
functional & \texttt{\nolinkurl{NO_NATIVE_PORT}} \\
Finite signed path pairing & Partial sign-reversing path involution &
\texttt{\nolinkurl{NO_NATIVE_PORT}} \\
\end{longtable}

The
\href{https://github.com/dbccpage/the_symbol_game/blob/52cb5457f6c485ba548cfd292f6012a941ed22cf/omega_engine_symbol_game/use_cases/mssn_sn_demos/demo_3/host.py}{use-case
host} supplies phases, fuel accounting, atomic transaction rollback, and
hash-chained batch receipts. Each batch that the host processes costs
one unit of fuel, charged before its moves are applied; the default
initial fuel is 12. If the batch is refused, the semantic state (covered
elements, witness candidate, phase) returns to its pre-batch value, but
the fuel unit is not returned. A batch with a move after
\passthrough{\lstinline!VERIFY!} or with a malformed element identifier
is refused and charged in this way. An empty batch is also charged; it
is recorded as \passthrough{\lstinline!ACCEPTED!} and changes no
semantic state. Once fuel is zero, the next batch is not applied: it is
recorded as \texttt{\nolinkurl{FUEL_EXHAUSTED}} and the host
enters \passthrough{\lstinline!TERMINAL!}, after which batches are
rejected with \texttt{\nolinkurl{REJECT_TERMINAL_STATE}}. The
number of processed batches is therefore bounded by the initial fuel.
This does not bound all attempted calls, receipt history, input bytes,
or wall-clock time. After \passthrough{\lstinline!TERMINAL!}, further
calls still pass through \texttt{\nolinkurl{_commit}} and append
refusal receipts at the baseline. Batch length and serialization size
also require caller limits; fuel bounds charged batches only.

The host's \texttt{\nolinkurl{state_digest}} starts from a SHA-256
digest of the fixture and arm identifiers. Each batch extends it with
SHA-256 over the prior digest, the proposed moves, and the host
decision. The moves are serialized as JSON with sorted keys, no white
space, and ASCII escapes; this is a declared convention, not a full RFC
8785 implementation. The digest does not cover fuel or the rest of the
host state. The receipt verifier
\texttt{\nolinkurl{verify_receipt_bytes}} accepts only the
canonical encoding of the expected 30-cell matrix (five fixtures × six
arms). For each cell it recomputes the digest chain from the initial
digest, checks the fuel sequence and the scripted decision sequence for
the arm, and refuses a receipt that reports a terminal certificate or a
witness type while the route port is missing. It therefore rejects
\passthrough{\lstinline!\{\}!}, schema changes, chain discontinuities,
and certificate claims. A SHA-256 chain without a key or an
independently held anchor shows internal consistency; it does not
authenticate who produced the receipt.

The host's \texttt{\nolinkurl{verify_route}} function returns a
typed missing-port refusal for any nonempty route name. The module does
not call a native board or quotient action. While the seam is missing,
returning \texttt{\nolinkurl{NO_NATIVE_PORT}} is the appropriate
outcome: it reports that verification is unavailable, not that a witness
is wrong.

\hypertarget{the-package-re-implementation}{%
\subsubsection{4.1 The package
re-implementation}\label{the-package-re-implementation}}

So that the refusal behavior can be run from this deposit alone, the
package includes a small re-implementation:
\texttt{\nolinkurl{code/host_engine.py}} (host),
\texttt{\nolinkurl{code/danger_gate.py}} (two-signal table),
\texttt{\nolinkurl{code/canaries.py}} (phrase detectors), and
\texttt{\nolinkurl{code/self_panel.py}}, with a Rust port of the
first three in \texttt{\nolinkurl{code/rust/}}. It is not the
canonical host, and its tests are illustrations, not integration
evidence. It differs from the baseline in these respects:

\begin{itemize}
\tightlist
\item
  its default fuel is 11, and it enters
  \passthrough{\lstinline!TERMINAL!} as soon as a batch spends the last
  unit of fuel;
\item
  its fixture list replaces
  \texttt{\nolinkurl{cochain_nonexact_cycle}} with
  \texttt{\nolinkurl{frustrated_square_diagonal}}, and it uses
  five route names, only one of which matches the baseline;
\item
  it refuses an empty batch
  (\texttt{\nolinkurl{REJECT_EMPTY_BATCH}}), and its
  \passthrough{\lstinline!VERIFY!} step requires only nonempty coverage
  and no witness;
\item
  it represents the Arm D detector hold as a move type, not as a host
  method that records detector findings;
\item
  its digests use a JSON initial record, the hexadecimal form of the
  move bytes, and UTF-8 output rather than ASCII escapes;
\item
  its receipt verifier recomputes the digest chain and refuses
  certificate labels, but it does not check fuel or the scripted arm
  sequences, so an edited receipt whose digests are recomputed passes
  unless it carries a certificate label;
\item
  its gate labels every two-signal refusal
  \passthrough{\lstinline!QUARANTINED!} and counts any supplied Signal 2
  record as present, whereas the baseline gate counts a record only when
  its defect count is positive;
\item
  \texttt{\nolinkurl{code/self_panel.py}} builds a synthetic
  1,000-row panel with the same shape (40 base strings, 25 tagged copies
  each) instead of loading the repository panel.
\end{itemize}

\hypertarget{worked-example-a-plausible-witness-that-cannot-be-certified}{%
\subsection{5. Worked example: a plausible witness that cannot be
certified}\label{worked-example-a-plausible-witness-that-cannot-be-certified}}

Initialize
\texttt{Demo3HostEngine("gauge\_even\_triangle", "T", initial\_fuel=12)}
in the baseline use-case host. The host begins in
\passthrough{\lstinline!INTAKE!}, with empty coverage and no witness.
Submit one batch containing:

\begin{lstlisting}
[
  {"move_type": "COVER_ELEMENT", "payload": {"element_id": "elem_all"}},
  {"move_type": "PROPOSE_WITNESS",
   "payload": {"witness_type": "vertex_bits", "witness": [0, 1, 0]}},
  {"move_type": "VERIFY", "payload": {}}
]
\end{lstlisting}

The example displays the proposal fields; the tests construct typed
\passthrough{\lstinline!MoveProposal!} objects. The host charges one
unit of fuel, records coverage and the witness provisionally, then
reaches \passthrough{\lstinline!VERIFY!}. Since the required route port
is absent, it returns \texttt{\nolinkurl{NO_NATIVE_PORT}} and
rolls back the batch. Coverage is empty again, the witness is absent,
the phase is \passthrough{\lstinline!INTAKE!}, and no terminal
certificate exists. Remaining fuel is eleven: the charge is made before
the moves are applied and is not returned on rollback. The step receipt
records the phases before and after the batch, the proposed moves, the
refusal decision, the prior and new digests, and the remaining fuel. The
package re-implementation gives the same outcome for this batch; its
test
\texttt{\nolinkurl{test_known_gauge_witness_still_refuses_and_rolls_back}}
checks it, including the remaining fuel.

Now alter the batch by appending another move after
\passthrough{\lstinline!VERIFY!}. The host refuses with
\texttt{\nolinkurl{REJECT_MOVE_AFTER_VERIFY}}, charges one unit,
and rolls back the whole batch. A batch with a malformed element
identifier is refused in the same way with
\texttt{\nolinkurl{REJECT_MALFORMED_PAYLOAD}}. An empty batch is
charged one unit but changes no semantic state; the package
re-implementation refuses it instead. After the twelfth charged batch,
remaining fuel is zero; the next batch is not applied, is recorded as
\texttt{\nolinkurl{FUEL_EXHAUSTED}}, and leaves the host in
\passthrough{\lstinline!TERMINAL!}.

For a novice, the distinction is practical: ``the checker is
unavailable'' is an answer about the verification process. It is not an
answer that the proposed mathematical witness is false. The retained
attempt can later be replayed against the same protocol; testing the
witness after a port is added requires a new, explicitly versioned
evaluation.

\hypertarget{design-for-a-prespecified-effectiveness-experiment}{%
\subsection{6. Design for a prespecified effectiveness
experiment}\label{design-for-a-prespecified-effectiveness-experiment}}

The experiment is conditional on implementing and independently
validating the three route adapters. The existing
\href{demos/demo_3/experiment_manifest.yaml}{manifest} remains
\texttt{\nolinkurl{proposed_not_executed}}, with live execution
disabled and provider/model fields unset. The following design must be
frozen before collecting confirmatory data.

Use a finite task bank stratified by route, difficulty, satisfiable
versus obstructed outcome, and vocabulary relabeling. Reserve task
families and a second model for held-out evaluation. Discovery examples,
admission examples, pilot examples, and confirmatory examples must be
disjoint at the underlying instance level. Surface renaming alone is
insufficient to make a new task.

The six planned arms are plain proposal (P), typed protocol (T),
structured coverage guidance (C), generic reflection (R), admitted
detector policy with feedback (D), and solution-bearing assistance (A).
Every arm uses the same output schema, route checker, stopping budget,
and task inputs. Arm A is an assisted control and is excluded from
claims about unassisted solution finding. The principal contrast is D
minus R: the combined detector policy and its feedback, relative to
generic reflection. Match feedback opportunities, prompt and token
accounting, and stopping budgets; freeze the differing feedback content.
This contrast estimates the combined policy, not the detector's isolated
mechanism.

Randomize arm order within fixture/model blocks and record the
assignment. Model revisions, provider route, temperature, seed support,
prompt templates, detector bank, checker bytes, retry policy, and
budgets must be pinned. Unsupported seeds are reported as unsupported;
repeated requests then remain stochastic trials. The pilot estimates
resource requirements and outcome variation. It does not determine
success thresholds after seeing the confirmatory result.

The primary effectiveness outcome is a run-level binary variable: an
independently verified route witness obtained within the common budget.
Unknown, refused, timed-out, and exhausted runs are retained in the
denominator and reported separately. Verification errors are not
silently converted into incorrect model answers. A checker defect
discovered during collection stops the affected analysis and requires a
new source-bound run.

The confirmatory analysis should estimate the D--R difference using
fixture blocks as the dependence unit. A prespecified paired or
randomization analysis must respect the actual allocation. Report route
strata, model strata, the difference, and its uncertainty. The
effectiveness gate is a strictly positive lower endpoint of a 95\%
interval for the principal held-out comparison. Exploratory comparisons
are labeled exploratory, with their multiplicity reported.

\hypertarget{safety-transfer-and-stopping-criteria}{%
\subsection{7. Safety, transfer, and stopping
criteria}\label{safety-transfer-and-stopping-criteria}}

For route \(r\), false suppression is the conditional probability of a
block under the frozen policy among independently valid,
detector-eligible proposals. Its observed numerator is valid eligible
blocked proposals and its denominator is all valid eligible proposals.
Blocked invalid proposals do not enlarge that denominator. In D, log
proposals blocked by the enforcement policy and evaluate them with the
same route oracle under a safe shadow path. In R, run the frozen
detector in shadow mode on observed proposals without changing feedback.
Record eligibility and whether the independent oracle accepts; do not
infer validity from a model's declaration.

For every route proposed for use, require an upper 95\% confidence bound
below 1\% on false suppression and no false terminal certificate. These
are individual route bounds; a simultaneous safety claim across routes
needs prespecified multiplicity handling. Fixed-sample confidence
calculations do not license optional repeated looks without an
appropriate sequential rule. A route without enough valid eligible
proposals is unassessed. Effectiveness on one route cannot substitute
for safety on another. Adversarial near misses supplement observed
proposal streams but must have separate denominators.

The sample-size issue is concrete. Under independent Bernoulli trials,
observing zero failures in n opportunities gives the one-sided exact
upper bound \texttt{1 - 0.05\^(1/n)}, obtained by
solving \texttt{(1-p)\^n = 0.05}. Eight zero-failure
observations leave a bound near 31\%; 299 independent opportunities are
needed to bring it below 1\%. Shared fixtures and correlated proposals
invalidate a naive independence assumption, so these numbers are an
illustration of the planning burden, not this study's achieved safety
result.

Held-out transfer evaluates the frozen detector bank without editing its
rules. Report failures caused by new task families and model behavior as
well as successes. If a new detector is added, restart the relevant
held-out assessment. Stop a run on its declared call, token, time, or
monetary budget. Stop the study on a false terminal, compromised source
binding, or loss of required custody records. Budget exhaustion remains
an observed outcome.

\hypertarget{evidence-custody-and-reproducibility}{%
\subsection{8. Evidence custody and
reproducibility}\label{evidence-custody-and-reproducibility}}

The historical pilot file in this package,
\texttt{\nolinkurl{demos/demo_3/results/live_openrouter_trial_receipt.json}},
contains descriptive arm counts and provider fields for one model and
two arms. Its raw request/source bindings do not establish the frozen
randomized experiment above, and its trial records report terminal
certificates from an earlier runner that the audited host cannot issue.
Consequently this paper reports no live effect estimate, confidence
interval, cost advantage, or transfer result from it.

A confirmatory receipt must bind the task, detector bank, checker,
assignment, model route, parameters, raw authorized proposal bytes,
decisions, state transitions, and measured resource use. Retention
requires the active data policy; refusing a proposal does not authorize
keeping secrets or personal data. A hash supports checking recorded
bytes against an independently held anchor. It does not authenticate a
provider or prove the semantic verdict.

\hypertarget{package-checks}{%
\subsubsection{8.1 Package checks}\label{package-checks}}

From the package root, with Python 3.12, install the pinned tools and
run:

\begin{lstlisting}
py -3.12 -m pip install "pytest>=9.0.3" ruff==0.14.13
py -3.12 -m pytest -o addopts= code/tests -q -p no:cacheprovider
py -3.12 -m ruff check code
cargo test -j 1 --manifest-path code/rust/Cargo.toml
\end{lstlisting}

The Rust-versus-Python comparison in
\texttt{\nolinkurl{code/tests/test_danger_gate.py}} runs only
when the binary has been built
(\texttt{cargo build --release -j 1 --manifest-path code/rust/Cargo.toml})
or its path is given in
\texttt{\nolinkurl{PFP003_HOST_DETECTOR_BIN}}; otherwise it is
skipped with that reason. These checks exercise the package
re-implementation of Section 4.1, not the canonical host.

\hypertarget{repository-checks-requires-the-full-repository-checkout}{%
\subsubsection{8.2 Repository checks (requires the full repository
checkout)}\label{repository-checks-requires-the-full-repository-checkout}}

The following commands require the full repository checkout at the
baseline commit; they are not runnable from this deposit. Collect the
focused tests before execution:

\begin{lstlisting}
python -m pytest omega_engine_symbol_game/tests/mssn_sn_demos/test_demo_3.py --collect-only -q
python -m pytest omega_engine_symbol_game/tests/mssn_sn_demos/test_demo_3.py -q
python -m pytest omega_engine_symbol_game/tests/test_immune_defect_regressions.py -q
\end{lstlisting}

The
\href{https://github.com/dbccpage/the_symbol_game/blob/52cb5457f6c485ba548cfd292f6012a941ed22cf/omega_engine_symbol_game/tests/mssn_sn_demos/test_demo_3.py}{Demo
3 tests} include the refusal, rollback, fuel, and matrix receipt checks
described in Sections 4 and 5. The two-signal conjunction example is the
test
\texttt{\nolinkurl{test_invariant_two_signal_danger_gate_conjunction_logic}}
in
\texttt{\nolinkurl{omega_engine_symbol_game/tests/mssn_sn_demos/test_demo_immune_soviet.py}}.
The catalog integrity suite is
\texttt{\nolinkurl{omega_engine_symbol_game/tests/test_mssn_catalog.py}}.
Its offline fixture blocks sockets and distinguishes reproducible
fixtures from live receipts whose verification is integrity-only. Use a
fresh writable temporary directory; record the interpreter, collection
status, command exit, and source digests. These commands are a
reproducibility plan; current results must be read from the associated
execution receipt.

\hypertarget{threat-model-and-limitations}{%
\subsection{9. Threat model and
limitations}\label{threat-model-and-limitations}}

An untrusted proposer may supply a verification flag, reuse a stale
witness, name nonexistent coverage, or exploit a missing adapter. An
erroneous detector may reject valid novelty. A defective host may accept
a malformed transition. The protocol addresses these threats through
typed inputs, negative examples, independent route outcomes, and
recorded refusals, each within its tested scope.

It does not assume that all host code is correct solely because it is
outside the model. The canonical route seam is missing at its audited
baseline; signal evidence is supplied to the filter; the self panel has
limited diversity; the four phrase detectors are unadmitted and cannot
see negation or context; and illustrative package tests are not
integration evidence. Fuel alone does not bound serialization,
post-terminal receipt retention, or wall-clock cost. Formal models of
fuel or authority apply only to their stated types and do not prove
correspondence to the Python host.

The resulting contribution is a concrete admission and evaluation method
with inspectable implementation boundaries. Whether admitted detectors
improve solution finding safely remains an empirical question with
explicit conditions for an answer. This paper assigns no authorization
and changes no jurisdiction. It carries no operational authority or
institutional standing.

\hypertarget{verification-for-this-revision}{%
\subsection{Verification for this
revision}\label{verification-for-this-revision}}

The package verifier (\texttt{\nolinkurl{verify_deposit.py}})
checks the manifest and the two-way file census, and
\texttt{tools/rebuild\_release.py --check} checks the
release archive against the manifests. The fresh package checks are
recorded in
\texttt{\nolinkurl{docs/first_publication_validation_2026_10_02.json}},
with command, tool version, exit status, input hashes, and skipped or
unavailable scope. The review is in
\texttt{\nolinkurl{docs/first_publication_review_2026_10_02.md}}.
The revised 24-problem Landau quiz, study guide, and separate instructor
key are in
\texttt{\nolinkurl{_lev_landau_theoretical_minimum_quiz/}}.
Historical receipts in this package keep their original dates and scope
and are not current execution evidence.

\hypertarget{appendix-a.-two-signals-one-detector}{%
\subsection{Appendix A. Two signals, one
detector}\label{appendix-a.-two-signals-one-detector}}

Section 3 describes two safeguards. A pattern match (Signal 1) should
refuse a proposal only when a structural defect (Signal 2) corroborates
it, and a candidate detector should be screened against a self panel of
strings treated as valid. Each safeguard rests on an assumption that is
easy to lose in an implementation. This appendix states three small
lemmas that make the assumptions exact, checks each one exhaustively on
finite domains, audits the repository against them, and describes one
component that realizes the admission record of Section 3.

The checks are in \texttt{\nolinkurl{code/detector_lemmas.py}} and
\texttt{\nolinkurl{code/tests/test_two_signals_one_detector.py}}.
They use exact integers and fractions and import nothing from the
repository. Match counts, rational signal inputs, and predicate/panel
masks refuse floats and Booleans. The exact threshold helper
additionally requires a positive integer exponent and positive scale;
out-of-universe predicate bits are refused by the universal panel
checks. Unless stated otherwise, repository paths in this appendix refer
to commit
\texttt{\nolinkurl{30609de2c3928238715375a9bac608256fa381d8}}, the
revision audited on 2026-09-28, not to the code baseline of the rest of
the paper.

\hypertarget{a.1-the-decision-table}{%
\subsubsection{A.1 The decision table}\label{a.1-the-decision-table}}

Write \(m\) for the number of Signal 1 matches and \(s\) for the Signal
2 record, which counts as present when its defect count is positive. The
table of Section 3, as implemented in the repository gate
(\texttt{\nolinkurl{omega_engine_epistemic_immune_system/domain/danger_gate.py}}),
gives \passthrough{\lstinline!TOLERATED!} when neither signal is present
and \texttt{\nolinkurl{DIAGNOSTIC_ONLY}} when only Signal 1 is.
When only Signal 2 is present it gives
\texttt{\nolinkurl{DEFENSIVE_LOCKDOWN}} for a high or critical
defect and \texttt{\nolinkurl{DIAGNOSTIC_ONLY}} otherwise. When
both are present it gives \passthrough{\lstinline!QUARANTINED!} for a
high or critical defect and
\texttt{\nolinkurl{DEFENSIVE_LOCKDOWN}} otherwise.
\texttt{\nolinkurl{DIAGNOSTIC_ONLY}} is the outcome the design
exists for: a match alone is recorded but refuses nothing.

\hypertarget{a.2-lemma-a-a-derived-second-signal-collapses-the-table}{%
\subsubsection{A.2 Lemma A: a derived second signal collapses the
table}\label{a.2-lemma-a-a-derived-second-signal-collapses-the-table}}

\textbf{Lemma A.} Suppose Signal 2 is computed from Signal 1,
\(s = g(m)\) for some function \(g\). Then:

\begin{enumerate}
\def\labelenumi{\arabic{enumi}.}
\tightlist
\item
  the decision is a function of \(m\) alone, so two proposals with the
  same number of matches receive the same decision whatever their
  structural defects;
\item
  \texttt{\nolinkurl{DIAGNOSTIC_ONLY}} is reachable exactly when,
  for some \(m \geq 1\), \(g(m)\) has no positive defect, or \(g(0)\) is
  a positive defect of low or medium severity;
\item
  if Signal 2 is present exactly when Signal 1 is, with one severity
  \(\sigma\) for every \(m \geq 1\), then exactly two decisions are
  reachable: \passthrough{\lstinline!TOLERATED!} and
  \passthrough{\lstinline!QUARANTINED!} when \(\sigma\) is high or
  critical, and \passthrough{\lstinline!TOLERATED!} and
  \texttt{\nolinkurl{DEFENSIVE_LOCKDOWN}} otherwise.
\end{enumerate}

\emph{Proof.} Part 1 is substitution. For part 2, the table yields
\texttt{\nolinkurl{DIAGNOSTIC_ONLY}} only in the cell where Signal
1 is present and Signal 2 is absent, and in the cell where Signal 2
alone is present with low or medium severity; a derived Signal 2 reaches
these cells exactly under the stated conditions. For part 3, the
hypothesis leaves only the cells where both signals are absent
(\(m = 0\)) and where both are present with severity \(\sigma\)
(\(m \geq 1\)), and both cells are reached. \(\square\)

Under part 3 with a high severity the gate refuses exactly when Signal 1
is present. It is then a one-signal detector, and a match can no longer
be recorded without refusing. The test
\texttt{\nolinkurl{test_derived_signal2_exhaustive_over_every_function_of_the_match_count}}
enumerates all \(13^4 = 28{,}561\) functions from the match counts
\(\{0, 1, 2, 3\}\) to the 13 Signal 2 values (none, or one of four
severities with 0, 1 or 2 defects). It confirms part 2 for every
function and part 3 for the 160 functions that satisfy its hypothesis.
With independently selectable test inputs all four decisions occur
(\texttt{\nolinkurl{test_independent_signals_reach_all_four_decisions}}).
This is table coverage, not proof of independent derivation in a
deployed pipeline. The package's own table of Section 4.1 collapses in
the same way
(\texttt{\nolinkurl{test_package_gate_collapses_the_same_way_under_a_derived_signal2}}).

\hypertarget{a.3-lemma-b-the-hill-type-quorum-is-a-sum-threshold}{%
\subsubsection{A.3 Lemma B: the Hill-type ``quorum'' is a sum
threshold}\label{a.3-lemma-b-the-hill-type-quorum-is-a-sum-threshold}}

The repository also has a gate described as a cooperative quorum
(\texttt{\nolinkurl{omega_engine_epistemic_immune_system/domain/allostery.py}}).
It adds the strictly positive sensor signals into
\(\theta = \sum_i s_i\), computes the Hill ratio
\(H(\theta) = \theta^n / (K^n + \theta^n)\) with \(H(0) = 0\), and
activates when \(H(\theta) \geq t\).

\textbf{Lemma B.} Let \(n > 0\) and \(K > 0\).

\begin{enumerate}
\def\labelenumi{\arabic{enumi}.}
\tightlist
\item
  For \(0 < t < 1\) the gate activates exactly when
  \(\theta \geq \theta^* = K \, (t / (1 - t))^{1/n}\).
\item
  At \(t = 1/2\), \(\theta^* = K\) for every \(n\), so the Hill
  coefficient does not affect the decision. For other \(t\), \(n\) only
  moves the threshold: the gate with parameters \((n, K, t)\) decides
  exactly as the gate with parameters \((1, \theta^*, 1/2)\).
\item
  The gate implies no number of agreeing sensors. One signal of size
  \(\theta^*\) activates it. Only if every signal is at most \(u\) does
  activation require at least \(\lceil \theta^* / u \rceil\) positive
  signals.
\item
  At \(t = 0\) it activates with no signal at all, and at \(t = 1\) it
  never activates.
\end{enumerate}

\emph{Proof.} For \(\theta > 0\), \(H(\theta) \geq t\) holds exactly
when \(\theta^n (1 - t) \geq t K^n\), that is when
\(\theta^n \geq K^n t / (1 - t)\); both sides are positive and
\(x \mapsto x^{1/n}\) is increasing, which gives part 1. At
\(\theta = 0\), \(H = 0 < t\). Part 2 follows because
\(t / (1 - t) = 1\) at \(t = 1/2\), and because part 1 describes both
gates by the same threshold. For part 3, \(m\) signals each at most
\(u\) sum to at most \(mu\). For part 4, \(H(0) = 0 \geq 0\), and
\(H(\theta) < 1\) for every finite \(\theta\). \(\square\)

The tests check the lemma with exact fractions for integer \(n\) from 1
to 4, four values of \(K\) and five of \(t\), over all 400 signal
vectors of length at most 3 drawn from seven values, including 0 and a
negative value: the ratio computed as a fraction agrees with the
cleared-denominator threshold, the decision depends only on \(\theta\)
and increases with it, a single sensor activates every threshold, and
the bounded-sensor count holds. At \(t = 1/2\) the decisions agree for
\(n = 1, \ldots, 6\). The repository's defaults are \(n = 5/2\),
\(K = 2\) and \(t = 1/2\); by part 2 the gate then activates when the
positive signals sum to at least 2, whatever \(n\) is. A per-sensor cap
yields a count of positive entries, not necessarily distinct or
independent agents; identity, deduplication, and provenance require
additional checks.

\hypertarget{a.4-lemma-c-negative-selection-screens-what-it-screens}{%
\subsubsection{A.4 Lemma C: negative selection screens what it
screens}\label{a.4-lemma-c-negative-selection-screens-what-it-screens}}

A detector rule has a \emph{screened} predicate \(P\), which the screen
tests against the self panel \(S\), and an \emph{executed} predicate
\(E\), which the matcher computes on the content it inspects. The screen
admits the rule when no \(x \in S\) satisfies \(P\). Screening is
\emph{sound} for \(S\) if an admitted rule never fires on a member of
\(S\).

\textbf{Lemma C.} Over a universe \(U\) of inputs:

\begin{enumerate}
\def\labelenumi{\arabic{enumi}.}
\tightlist
\item
  screening is sound for every finite panel \(S \subseteq U\) exactly
  when \(E(x)\) implies \(P(x)\) for every \(x \in U\);
\item
  for every finite panel, screening never rejects a rule that would not
  fire on that panel exactly when \(P(x)\) implies \(E(x)\) for every
  \(x \in U\);
\item
  both hold for every panel exactly when \(E = P\) on \(U\).
\end{enumerate}

\emph{Proof.} For part 1, if \(E\) implies \(P\) and the rule is
admitted, no member of \(S\) satisfies \(P\), so none satisfies \(E\).
Conversely, if some \(x\) has \(E(x)\) but not \(P(x)\), the panel
\(S = \{x\}\) admits the rule, which then fires on \(x\). Part 2 is the
same argument with \(E\) and \(P\) exchanged, and part 3 combines the
two. \(\square\)

Even \(E=P\) guarantees only the stated panel-relative property. It does
not show that the panel contains every valid input, nor estimate a
false-suppression rate outside it.

The test
\texttt{\nolinkurl{test_screening_is_sound_for_every_panel_iff_executed_implies_screened}}
checks parts 1 and 2 for all \(64 \times 64\) pairs of predicates on a
six-element universe and all 63 nonempty panels. The test
\texttt{\nolinkurl{test_pattern_ignoring_matcher_is_admitted_and_fires_on_a_self_fixture}}
gives the counterexample of part 1 in the form found below: the screen
tests a failure-class name, the matcher tests the structure of the
record, and the rule is admitted and fires on a member of the panel.

\hypertarget{a.5-the-repository-audited-against-the-lemmas}{%
\subsubsection{A.5 The repository audited against the
lemmas}\label{a.5-the-repository-audited-against-the-lemmas}}

The findings below refer to
\texttt{\nolinkurl{omega_engine_epistemic_immune_system/}} at
the audited commit. Correction paragraphs refer to the separate
historical revision
\texttt{\nolinkurl{fdc0035a81c45e10b900fc917b8b0710df81391f}}; they
do not describe the default branch or a fresh integration test. That
correction changes the gate policy as well as signal provenance.

\emph{Lemma A.} The packaged demonstration
\texttt{\nolinkurl{use_cases/demos/demo_4e_cogsci_autopoiesis/run_packaged_immune.py}}
computes Signal 2 for each of its 24 control cases from the match count:
line 178 passes
\texttt{\{"defect\_count": len(matches)\}} to a sensor
that reports a high-severity defect whenever the count is positive
(\texttt{\nolinkurl{infrastructure/mssn_ports.py}}, lines 27 to
39). Signal 1 is then recomputed by the same matcher over the rules that
the screen admits. When the screen admits every rule, part 3 of Lemma A
applies: the only reachable decisions are
\passthrough{\lstinline!TOLERATED!} and
\passthrough{\lstinline!QUARANTINED!}, and every match refuses. The
demonstration's other path, for model replies, takes Signal 2 from a
separate defect check and is not affected. Correction: in commit
\passthrough{\lstinline!fdc0035a8!} on branch
\texttt{\nolinkurl{claude/mssn-study-corrections}}, Signal 2 comes
from host checks that the matcher never sees
(\texttt{\nolinkurl{domain/host_checks.py}} and
\texttt{\nolinkurl{domain/danger_signals.py}}),
\texttt{\nolinkurl{assert_signals_independent}} refuses a derived
Signal 2 with the code \texttt{\nolinkurl{TWO_SIGNAL_COLLAPSE}},
and the gate follows the two-signal rule, under which Signal 2 alone is
diagnostic only
(\texttt{\nolinkurl{tests/unit/test_two_signal_independence.py}}).

\emph{Lemma B.} \texttt{\nolinkurl{domain/allostery.py}}, lines 74
to 79, computes the gate of Section A.3 in floating point. At its
defaults it activates when the positive signals sum to at least 2, and a
single sensor suffices. The comment at line 68 says that \(K = 2\)
``requires 2 equivalent unit signals''; that holds only if each signal
is at most 1, which the code does not require. With \(t = 0\) it
activates on an empty list. Correction: in commit
\passthrough{\lstinline!fdc0035a8!} the gate is exact, with an integer
Hill coefficient, a threshold strictly between 0 and 1, a cap of 1 on
each signal, so that at the defaults at least two sensors are needed,
and typed refusals
(\texttt{\nolinkurl{tests/unit/test_allostery_exact.py}}).

\emph{Lemma C.} The screen in
\texttt{\nolinkurl{domain/thymus.py}}, lines 77 to 99, compiles
each rule's regular expression and searches the self panel with it.
Three of the four matchers behind the same port
(\passthrough{\lstinline!PharmacyMatcher!},
\passthrough{\lstinline!FarmMatcher!} and
\passthrough{\lstinline!ProtocolMatcher!} in
\texttt{\nolinkurl{infrastructure/mssn_ports.py}}, lines 80 to
166) never execute that expression: they use a rule only to select a
target class and report the findings of a separate detector. For the
pharmacy rules the screened expression is the escaped name of the
failure class (\texttt{\nolinkurl{run_packaged_immune.py}}, lines
118 to 122). The executed and screened predicates differ, so by Lemma C
the screen gives these rules no soundness: a self fixture that the
detector flags but that does not contain the class name passes the
screen. Correction: in commit \passthrough{\lstinline!fdc0035a8!} each
rule binds the identifier of the matcher that executes it, the screen
runs that matcher on the self panel, and regular expressions are limited
to a subset that cannot backtrack without bound
(\texttt{\nolinkurl{domain/safe_pattern.py}};
\texttt{\nolinkurl{tests/unit/test_thymus_binds_executed_matcher.py}}).

The package re-implementation of Section 4.1 is not affected by the last
two findings; its table counts any supplied Signal 2 record as present,
and Lemma A applies to it whenever that record is computed from the
matches.

\hypertarget{a.6-a-realized-admission-record-anti-skill-v2}{%
\subsubsection{A.6 A realized admission record: anti-skill
V2}\label{a.6-a-realized-admission-record-anti-skill-v2}}

Section 3 asks that a detector be admitted only with a reproduced
defect, valid near misses, a frozen predicate and bound digests. One
component of the repository states that record as types: anti-skill V2,
in
\texttt{\nolinkurl{omega_engine_workflows_agents_skills/domain/anti_skill/anti_skill_v2.py}}
and \texttt{\nolinkurl{application/anti_skill/anti_skill_v2.py}}
of the same package.

\begin{itemize}
\tightlist
\item
  \emph{A closed predicate language.} A detector is an expression of at
  most 16 atoms, combined by \passthrough{\lstinline!ALL!} or
  \passthrough{\lstinline!ANY!}. Each atom applies one of four operators
  (\texttt{\nolinkurl{EQUAL_FIELDS}},
  \texttt{\nolinkurl{NOT_SUBSET_FIELDS}},
  \texttt{\nolinkurl{VALUE_IN}},
  \texttt{\nolinkurl{IS_EMPTY}}) to eight declared fields of a
  subject record. There is no free-form code, regular expression or
  model verdict, so the predicate that is evaluated is the predicate
  that is recorded: the condition \(E = P\) of Lemma C holds by
  construction.
\item
  \emph{Positive and negative controls in at least two contexts.}
  Evaluation refuses a specification whose fixture bank lacks an
  expected hit, lacks a negative control, or draws its cases from fewer
  than two contexts (application module, lines 140 to 148). A negative
  control cannot expect a hit, and copied subjects are refused.
\item
  \emph{Derived receipts.} An evaluation receipt is refused unless its
  outcome equals the outcome computed from its case results, and every
  record's identifier is recomputed from its canonical content.
\item
  \emph{A human decision.} Admission requires a closed
  \passthrough{\lstinline!APPROVE!} or \passthrough{\lstinline!REJECT!}
  decision bound to the specification digest and to the evaluation
  receipt, and a verified registry requires one approving decision per
  specification, matched one to one with a passing receipt.
\end{itemize}

At the audited commit three specifications, their receipts and three
approving decisions exist in a registry scoped to one engine profile
(claim ledger row P3-26 gives the paths). The Python admission verifier
and the loader that would activate that registry are retired and return
a typed refusal
(\texttt{\nolinkurl{omega_engine_workflows_agents_skills/application/anti_skill_admission/verifier.py}},
lines 25 to 45), no Rust code reads the registry, and the registry that
the prompt reviewer loads is empty. At that audited commit, no detector
is therefore admitted and active through this record. Later integration
requires its own source audit and execution. It shows that the admission
record of Section 3 can be written as types whose outcomes are derived
rather than declared; whether admitted detectors help remains the
experiment of Sections 6 and 7. The corrections of 2026-09-29 leave
anti-skill V2 unchanged; moving it into the immune-system package, where
the corrected gate and screen now live, is recorded as open work.

\hypertarget{references}{%
\subsection{References}\label{references}}

\begin{itemize}
\tightlist
\item
  Chandrasekharan, S., and Wiese, U.-J. (1999). Meron-cluster solution
  of fermion sign problems. \emph{Physical Review Letters} 83,
  3116--3119.
  \href{https://doi.org/10.1103/PhysRevLett.83.3116}{doi:10.1103/PhysRevLett.83.3116};
  \href{https://arxiv.org/abs/cond-mat/9902128}{arXiv:cond-mat/9902128}.
\item
  Li, Z.-X., Jiang, Y.-F., and Yao, H. (2016). Majorana-time-reversal
  symmetries: a fundamental principle for sign-problem-free quantum
  Monte Carlo simulations. \emph{Physical Review Letters} 117, 267002.
  \href{https://doi.org/10.1103/PhysRevLett.117.267002}{doi:10.1103/PhysRevLett.117.267002};
  \href{https://arxiv.org/abs/1601.05780}{arXiv:1601.05780}.
\item
  Madaan, A., Tandon, N., Gupta, P., Hallinan, S., Gao, L., Wiegreffe,
  S., Alon, U., Dziri, N., Prabhumoye, S., Yang, Y., Gupta, S.,
  Majumder, B. P., Hermann, K., Welleck, S., Yazdanbakhsh, A., and
  Clark, P. (2023).
  \href{https://arxiv.org/abs/2303.17651v2}{Self-Refine: Iterative
  Refinement with Self-Feedback}. arXiv preprint 2303.17651, version 2.
\item
  Shinn, N., Cassano, F., Berman, E., Gopinath, A., Narasimhan, K., and
  Yao, S. (2023). \href{https://arxiv.org/abs/2303.11366v4}{Reflexion:
  Language Agents with Verbal Reinforcement Learning}. arXiv preprint
  2303.11366, version 4.
\item
  Troyer, M., and Wiese, U.-J. (2005). Computational complexity and
  fundamental limitations to fermionic quantum Monte Carlo simulations.
  \emph{Physical Review Letters} 94, 170201.
  \href{https://doi.org/10.1103/PhysRevLett.94.170201}{doi:10.1103/PhysRevLett.94.170201};
  \href{https://arxiv.org/abs/cond-mat/0408370}{arXiv:cond-mat/0408370}.
\item
  Wu, C., and Zhang, S.-C. (2005). Sufficient condition for absence of
  the sign problem in the fermionic quantum Monte Carlo algorithm.
  \emph{Physical Review B} 71, 155115.
  \href{https://doi.org/10.1103/PhysRevB.71.155115}{doi:10.1103/PhysRevB.71.155115};
  \href{https://arxiv.org/abs/cond-mat/0407272v3}{arXiv:cond-mat/0407272v3}.
\item
  Wei, Z. C., Wu, C., Li, Y., Zhang, S., and Xiang, T. (2016). Majorana
  positivity and the fermion sign problem of quantum Monte Carlo
  simulations. \emph{Physical Review Letters} 116, 250601.
  \href{https://doi.org/10.1103/PhysRevLett.116.250601}{doi:10.1103/PhysRevLett.116.250601};
  \href{https://arxiv.org/abs/1601.01994}{arXiv:1601.01994}.
\end{itemize}

Implementation links throughout the paper identify files at the stated
source baseline. The proposed experiment manifest is a separate design
artifact.

\end{document}
